\section{Ableitung des Forschungsfokus und Annahmen}
\label{sec:03-07_research-focus-and-assumptions}

Die Voranalyse zeigt, dass die zentrale Herausforderung der Plattform weniger in ihrer technischen Leistungsfähigkeit liegt als in der fehlenden systematischen Verzahnung von \Gls{ux}, \Gls{ia} und organisationalem Nutzungskontext. Im Vordergrund stehen daher nicht technische Einschränkungen, sondern die wahrgenommene Nützlichkeit und der Bedienungsaufwand des Systems. Diese gelten in Akzeptanzmodellen wie dem \acrshort{tam} und der \acrshort{utaut} als entscheidende Faktoren für die Nutzungsbereitschaft \cite{102:Davis1989,117:Horton2001,103:Venkatesh2003}. Vor diesem Hintergrund untersucht diese Arbeit empirisch, wie nutzerzentrierte \Gls{uxstrategien} und strukturelle Gestaltungsprinzipien die Orientierung, Akzeptanz und tatsächliche Nutzung interner Wissensplattformen beeinflussen.

\newpage

Aufbauend auf den im Praxisprojekt beobachteten Defiziten und Potenzialen wird angenommen, dass insbesondere drei Faktoren entscheidend für eine nachhaltige Nutzung sind:

\begin{tabular}{@{}p{1.2em}p{0.9\linewidth}@{}}
(1) & eine auf wahrgenommene Nützlichkeit ausgerichtete \Gls{ux}, die einen klar erkennbaren Beitrag zur Aufgabenerfüllung und Prozessunterstützung leistet (\Gls{performanceexpectancy}),\\[4pt]
(2) & die konsequente Berücksichtigung mentaler Modelle und Aufgabensichten der Nutzenden bei der Gestaltung von \Gls{ia} und Navigation sowie\\[4pt]
(3) & begleitende kommunikative Maßnahmen, die Sinn, Funktionsweise und Veränderungen der Plattform transparent machen \cite{103:Venkatesh2003,113:NNGIntranetGuidelines,118:NNGUserInvolvement}.
\end{tabular}


